\documentclass[10pt,a4paper]{amsart}
\usepackage[margin=19mm]{geometry}
\usepackage{amsmath,amssymb,mathtools}
\usepackage[scr=rsfs,cal=euler]{mathalfa}
\usepackage{xfrac}
\usepackage[unicode,hidelinks,bookmarks=false]{hyperref}
\newcommand{\ie}{i.e.\ }
\newcommand{\cf}{cf.\ }
\newcommand{\resp}{respectively\ }
\newcommand{\Subsection}{\subsection}
\providecommand{\nobreakdash}{\nobreak\hbox{-}}
\setlength{\emergencystretch}{2em}
\multlinegap0pt
\usepackage{amssymb}
\usepackage{mathtools}
\usepackage{tikz}
\newtheorem{claim}{Claim}
\title{Strong Majority Edge-Coloring}
\author{Sylwia Antoniuk, Magdalena Prorok, Nika Salia}
\date{Source: arXiv:2607.00212v1 (CC BY 4.0)}
\begin{document}
\maketitle

\noindent\emph{Source and licence.} Strong Majority Edge-Coloring. Version arXiv:2607.00212v1 (CC BY 4.0), distributed by arXiv under the Creative Commons Attribution 4.0 International licence, http://creativecommons.org/licenses/by/4.0/, as recorded in the arXiv licence metadata for this exact version and shown on the abstract page https://arxiv.org/abs/2607.00212v1. Full licence notice: \texttt{licenses/arxiv-2607.00212-CC-BY-4.0.txt}.\par\smallskip

\noindent\textit{Editorial scope.} Counted unit: the article's claim claim:extend (source0.tex lines 285--293). The article's complete original proof of it is delivered with it (source0.tex lines 295--313, one \texttt{proof} environment of 19 lines). No mathematics is written, repaired or re-derived: the statement and the proof are the article's own lines, in the article's own notation, and no retained line is substituted or re-flowed.  Retained uncounted as the source-local context the proof needs: lines 213--216; lines 225--227.  Nothing was omitted from any retained span, so the package carries no omission record.  No retained line contains a citation, so the wrapper carries no bibliography and no bibliography entry.  No retained line contains a figure environment, an image-inclusion command or drawing code, so no image is delivered and the package declares no asset dependency.  Editorial formatting only, in addition: an AMS \texttt{amsart} preamble that restates the article's own declarations and macros used by the retained lines, namely the environments \texttt{claim} on separate counters, together with the standard packages those lines need.  The retained environments print in this document as Claim 1 in order of appearance, whereas the article prints the same items with the numbering of its own full document.  The complete native source of the article as distributed in the arXiv e-print for this exact version is bundled unchanged as \texttt{sources/204154/source0.tex}.
\begin{equation}\label{eq:balanced}
        d_i(v)\;\le\;\frac{d_G(v)-1}{2}
        \qquad\text{for every color }i \in [k].
\end{equation}

\begin{claim}\label{claim:bad}
Let $d_G(v)=k\ge 3$ and suppose that \eqref{eq:balanced} holds at $v$. Call a color \emph{bad} for an edge $vw$ incident with $v$ if recoloring $vw$ with that color would violate \eqref{eq:balanced} at $v$. Then $vw$ has at most three bad colors if $k=4$, and at most two bad colors otherwise.
\end{claim}

\begin{claim}\label{claim:extend}
Let $H$ be a graph consisting of
\begin{itemize}
    \item a path $B=v\,e_0\,w_1\,e_1\cdots w_k\,e_k\,u$, where $k\ge 1$, $v\neq u$, and $uv\in E(G)$ whenever $k=1$,
    \item at least two edges incident with $v$ and different from $e_0$ (possibly with $uv$),
    \item at least two edges incident with $u$ and different from $e_k$ (possibly with $uv$).
\end{itemize} 
Suppose $d_H(v)\neq 5$ and $d_H(u)\neq 5$, and we have a coloring $c: E(H)\to [5]$, which is balanced at both $v$ and $u$, and which satisfies $c(e_0)\neq c(e_k)$ in case $k\geq 2$. Then $e_0,\dots,e_k$ can be recolored so that every edge of $B$ is strong majority colored and the new coloring is still balanced at both $v$ and $u$.
\end{claim}

\begin{proof}
We consider two cases.

\smallskip
\emph{Case $k=1$}
\smallskip

Note that we have $uv\in E(H)$. If $e_0$ and $e_1$ are already strong majority colored, we are done. If $d(v)\ge 4$ then the balanced condition at $v$ makes the edge $e_0$ strong majority colored, and symmetrically $e_1$ if $d(u)\ge 4$. Hence, an adjustment is needed only for the case $d_H(v)=3$ or $d_H(u)=3$. Suppose $d_H(v)=3$ and $e_0$ is not strong majority colored. We will recolor $e_1$. By Claim~\ref{claim:bad}, at most three colors are forbidden at $u$, and since $uv\in E(H)$, the color $c(uv)$ is one of them if there are exactly three colors forbidden. For the strong majority condition for the edge $e_0$, there are two colors forbidden, and one of them is $c(uv)$ again. Hence, a color remains for $e_1$ that keeps the balance at $u$ and makes $e_0$ strong majority colored. The symmetric failure at $u$ can be repaired in the same way by recoloring $e_0$. Note that by this procedure, no edge other than $e_0$ or $e_1$ will change its color.

Note that we have $uv\in E(H)$. If $e_0$ and $e_1$ are already strong majority colored, we are done. So suppose $e_0$ is not strong majority colored. We will adjust the color of $e_1$, preserving the balance at $u$. By Claim~\ref{claim:bad}, at most three colors are forbidden at $u$, and since $uv\in E(H)$, the color $c(uv)$ is one of them if there are exactly three colors forbidden. For the strong majority condition for the edge $e_0$, there are at most two colors forbidden for $e_1$, and in case of exactly two forbidden colors, one of them is $c(uv)$ again. Hence, a color remains for $e_1$ that keeps the balance at $u$ and makes $e_0$ strong majority colored. The symmetric failure at $e_1$ can be repaired in the same way by recoloring $e_0$. Note that by this procedure, no edge other than $e_0$ or $e_1$ will change its color.

\smallskip
\emph{Case $k\ge 2$.}
\smallskip

When $k=2$, the single interior edge $e_1$ is adjacent only to $e_0$ and $e_2$, which already have distinct colors, so $e_1$ satisfies the strong majority condition, regardless of its color. Now, since the coloring is balanced at $v$, there are at most two colors that, after being used to color $e_1$, could violate the strong majority condition for $e_0$, and the same holds for $e_2$. Therefore, there is a color available for $e_1$, i.e.~a color that makes both $e_0$ and $e_2$ strong majority colored. 

If $k>2$, we color the edges $e_j$, $j\in[k-1]$, one-by-one. Note that any such edge joins two degree-two vertices. Hence, while coloring $e_j$, we are restricted only by the strong majority condition for edges $e_{j-1}$ and $e_{j+1}$, and from each side there are no more than two forbidden colors. A free color for $e_j$ thus remains, and the coloring of the interior extends greedily. 
\end{proof}

\end{document}
